\documentclass[11pt,a4paper]{article}

\usepackage[utf8]{inputenc}
\usepackage[T1]{fontenc}
\usepackage{amsmath,amssymb,amsthm}
\usepackage{booktabs}
\usepackage{algorithm}
\usepackage{algpseudocode}
\usepackage[margin=2.6cm]{geometry}
\usepackage[hidelinks]{hyperref}

\newtheorem{proposition}{Proposition}
\newtheorem{remark}{Remark}

\newcommand{\Z}{Z}
\newcommand{\Dset}{D}
\newcommand{\E}{E(P_{\Dset})}
\newcommand{\R}{\mathbb{R}}
\newcommand{\Zint}{\mathbb{Z}}

\title{Hybrid methods for optimizing a linear fractional function\\
over the integer efficient set:\\
\large what each hybridization is worth, measured}

% Put your name and affiliation here before sending this anywhere.
\author{}
\date{}

\begin{document}
\maketitle

\begin{abstract}
Optimizing a linear fractional function $\Phi$ over the efficient set of a
multi-objective integer program is usually done by repeatedly truncating the
\emph{decision} space with Sylva--Crema cuts. Each cut states a disjunction, so
it costs $p$ binary variables and $p+1$ big-$M$ rows, and the sub-problem
solved at every iteration is strictly larger than the one before it. We first
move the unexplored part of the problem into \emph{criterion} space, as a list
of boxes: the truncation becomes a box subtraction, every sub-problem is the
original feasible set plus a handful of ordinary linear rows, and the model
never grows. On 35 instances this is faster on 34, by $5.05\times$ overall and
$9.01\times$ on the heaviest, with the margin widening in $p$.

We then hybridize that exact search three ways and report what each is worth.
Handing it the unfinished work of the published method (exact--exact)
\emph{ties} it. Seeding it from a Pareto local search with an archive
(exact--metaheuristic) gives $1.25\times$ against a measured ceiling of
$1.42\times$. Seeding it from the augmented weighted Tchebychev program --- the
scalarization the surrounding literature generates efficient points with ---
\emph{loses} to not seeding at all, at $0.71\times$.

A fourth change, pruning the open list the moment its best bound cannot beat
the incumbent, removes a quarter of the sub-problems and no measurable time;
stacked on the metaheuristic seed it removes 4\% rather than 25\%, because the
two feed the same bound comparison and are substitutes rather than
complements.

The Tchebychev result is the informative one, and it does not say the
scalarization is weak. As a generator it is strictly stronger than a weighted sum, reaching
unsupported efficient points no positive weighted sum can reach. It is a poor
\emph{seed} because it steers in criterion space relative to the ideal point
and never looks at $\Phi$. Symmetrically, the same free optimum that saves the
box search a third of its work saves the decision-space method \emph{zero}
iterations. What a heuristic incumbent is worth is a property of the exact
method it is attached to, not of the heuristic; we measure this in both
directions. All arithmetic is exact and rational, every answer is proved
optimal and cross-checked against independent references, and the negative
results and two corrected measurement errors are reported alongside the wins.
\end{abstract}

\tableofcontents

\section{Problem}

Let
\begin{equation}
  (P_{\Dset}) \qquad \text{``}\max\text{''} \quad
  \Z_k(x) = \frac{c_k^{\top}x + a_k}{d_k^{\top}x + b_k},
  \quad k = 1,\dots,p,
  \qquad x \in \Dset = \{\, x \in \Zint^n_+ : Ax \leqslant b \,\},
\end{equation}
with $\Dset$ non-empty and bounded and every denominator
$D_k(x) = d_k^{\top}x + b_k$ strictly positive on $\Dset$. A point
$x \in \Dset$ is \emph{efficient} when no $y \in \Dset$ satisfies
$\Z(y) \geqslant \Z(x)$ with at least one strict inequality; $\E$ denotes the
set of efficient points. Linear criteria are the degenerate case $d_k = 0$,
$b_k = 1$.

Given the decision maker's criterion
$\Phi(x) = (U^{\top}x + \alpha)/(V^{\top}x + \beta)$, the problem is
\begin{equation}
  (P_E) \qquad \max \ \Phi(x) \quad \text{s.t.} \quad x \in \E .
\end{equation}
$\E$ is a union of faces of $\Dset$: non-convex, with no explicit description.
The point of any method for $(P_E)$ is to reach the global optimum
\emph{without enumerating} $\E$.

An efficient point is \emph{supported} when it maximizes some strictly positive
weighted sum $w^{\top}\Z$ over $\Dset$ --- equivalently, when its criterion
vector lies on the convex hull of the criterion image. The distinction matters
throughout Section~\ref{sec:tcheb}: in an integer program unsupported efficient
points are common, and no weighted sum reaches any of them.

\section{Where the cost of the decision-space method comes from}

One iteration of the published method maximizes $\Phi$ over a truncated region
$\Dset_l$, tests the maximizer $x_l$ for efficiency, banks the best $\Phi$ on
the criterion slice about to be destroyed, and then truncates:
\begin{equation}\label{eq:cut}
  \text{delete} \quad \{\, x : \Z(x) \leqslant \Z(\tilde{x}) \,\} .
\end{equation}
Keeping the complement means keeping the $x$ for which \emph{some} criterion
strictly improves on $\tilde{x}$. That is a disjunction, and a disjunction is
not a system of linear inequalities. It is modelled with one binary per
criterion,
\begin{equation}\label{eq:bigM}
  e_k(x;\tilde{x}) \geqslant 1 \cdot y_k + M_k (1 - y_k), \quad k = 1,\dots,p,
  \qquad \sum_{k=1}^{p} y_k \geqslant 1 ,
\end{equation}
where $M_k$ is a lower bound on $e_k$ over the region, so $y_k = 0$ leaves its
row vacuous and $y_k = 1$ forces a strict improvement.

\textbf{Every cut therefore adds $p$ binaries and $p+1$ rows.} On the heaviest
instance we ship, eleven cuts leave $33$ binaries and $44$ rows on top of
$\Dset$, and the late sub-problems cost far more than the early ones. This is
the whole of the method's cost profile: the number of iterations is modest, but
each one is solved on a strictly larger mixed-integer program than the last.

This observation fixes the cost model used for every comparison in this note:
\begin{equation}\label{eq:cost}
  \text{total cost} \;\approx\; (\text{number of sub-problems solved})
  \times (\text{size of a sub-problem}).
\end{equation}
The two factors are attacked by different ideas. Section~\ref{sec:box} removes
the growth of the second. Sections~\ref{sec:meta} and~\ref{sec:tcheb} attack
the first, and the measurements there are only interpretable because the second
has already been made constant.

\section{The same truncation in criterion space}\label{sec:box}

Equation~\eqref{eq:cut} deletes a \emph{box} from $\R^p$. If the unexplored
part of the problem is carried as a list of boxes rather than as a region of
$\Zint^n$, the truncation is a rewrite of that list, and each sub-problem is
\begin{equation}\label{eq:sub}
  \max \ \Phi(x) \qquad \text{s.t.} \qquad x \in \Dset, \quad \text{the box's rows},
\end{equation}
which is $\Dset$ plus at most $2p$ ordinary linear rows. No binary variable
appears, and \textbf{the model never grows}: a box deep in the search carries a
few more rows than a shallow one, and that is all.

\subsection{How a box is written, and why not with numbers}

The obvious encoding, $l_k \leqslant \Z_k(x) \leqslant u_k$, fails. In the
general case $\Z_k$ is a ratio, so its values are rational, and the exact
``$+1$'' step that the split below needs would be lost.

Write the box with the linear form of the paper's Theorem~4 instead:
\begin{equation}\label{eq:e}
  e_k(x;a) \;=\; D_k(a)\, D_k(x) \bigl( \Z_k(x) - \Z_k(a) \bigr)
  \;=\; \bigl(D_k(a) c_k - N_k(a) d_k\bigr)^{\top} x
      + \bigl(D_k(a) a_k - N_k(a) b_k\bigr).
\end{equation}
Both denominators being positive, $e_k$ carries the sign of
$\Z_k(x) - \Z_k(a)$ while being \emph{linear in $x$}; with integer data it is
integer-valued on integer points. Hence, exactly and with no minimal step to
estimate,
\begin{equation}\label{eq:exact}
  \Z_k(x) > \Z_k(a) \iff e_k(x;a) \geqslant 1,
  \qquad
  \Z_k(x) \leqslant \Z_k(a) \iff e_k(x;a) \leqslant 0 .
\end{equation}
Linear and fractional criteria therefore go through the same code, and on a
linear criterion $e_k$ collapses to $c_k^{\top}x - c_k^{\top}a$.

\subsection{The split}

\begin{proposition}\label{prop:split}
Let $B$ be a box and $a \in \Dset$. For $k = 1,\dots,p$ put
\begin{equation}\label{eq:child}
  B_k \;=\; B \;\cap\; \{\, x : e_k(x;a) \geqslant 1 \,\}
            \;\cap\; \bigcap_{j<k} \{\, x : e_j(x;a) \leqslant 0 \,\} .
\end{equation}
Then $B_1,\dots,B_p$ are pairwise disjoint and
$\bigcup_k B_k = B \setminus \{\, x : \Z(x) \leqslant \Z(a) \,\}$.
\end{proposition}

\begin{proof}
Let $x \in B$ with $\Z(x) \not\leqslant \Z(a)$. The set of indices at which $x$
strictly beats $\Z(a)$ is non-empty; let $k$ be its smallest element. By
\eqref{eq:exact}, $e_k(x;a) \geqslant 1$, and $e_j(x;a) \leqslant 0$ for every
$j < k$, so $x \in B_k$. Conversely a point of $B_k$ has
$e_k \geqslant 1$, hence $\Z(x) \not\leqslant \Z(a)$. Disjointness: membership
of $B_k$ forces $k$ to be exactly that smallest index, and a point has only
one.
\end{proof}

\noindent
Each of the $p$ children is the parent's row list plus at most $k$ further
rows, all of them ordinary linear inequalities in $x$.

\subsection{Bounding, and why it is what makes seeding possible}
\label{sec:bound}

The sub-problem~\eqref{eq:sub} maximizes $\Phi$ over a \emph{superset} of the
efficient points whose criterion vector lies in the box, so its value bounds
all of them. A box whose bound does not beat the incumbent is therefore
discarded \emph{whole} and never split. Children inherit their parent's value
as a bound, which makes the open list a certified upper bound on everything
still unfound; the method is \emph{anytime} in the same sense, and for the same
reason, as the decision-space one.

This is the single structural fact that Sections~\ref{sec:meta}
and~\ref{sec:tcheb} rest on. A better incumbent does not merely stop the search
sooner; it deletes subtrees. The decision-space method has no such mechanism,
which is measured in Section~\ref{sec:asymmetry} and is the reason the two
methods respond so differently to the very same incumbent.

\subsection{Stopping when the best bound still open is hopeless}
\label{sec:exit}

The open list is kept in a heap ordered by inherited bound \emph{descending},
and only the root --- and the boxes a hybrid hands over --- carry no bound;
those rank first. So by the time a bounded box reaches the top, every box still
open is bounded too, and if this one cannot beat the incumbent, neither can any
of them. The search is finished, and the remaining list is busywork.

Measured before the change was made, on the 18 instances used throughout
Sections~\ref{sec:meta}--\ref{sec:tcheb}: 198 of 733 sub-problems (27\%) were
solved despite carrying an inherited bound the incumbent already beat --- and
on \emph{every single instance} that count equalled the number of boxes still
open when the condition first fired. The waste is exactly a tail, never
scattered, which is what the heap order predicts and which is why a single
early exit is the whole of the fix rather than a per-box filter. Stopping there
takes 733 boxes to 549, \textbf{25\% fewer}, with the same answers, all still
proved optimal.

\paragraph{It is not a speed-up, and we do not report it as one.} The total is
$7.44$\,s against $7.27$\,s at per-instance minima over five repeats, with
per-instance ratios scattered from $0.77\times$ to $1.36\times$. The change can
only \emph{remove} work, so every ratio below $1.00\times$ is machine noise ---
and the noise is larger than the effect.

That is a confirmation rather than a disappointment, and it sharpens the cost
model of \eqref{eq:cost}. The tail boxes are the \emph{cheapest} ones: each is
killed by the root relaxation the instant it is handed the cutoff, so a quarter
of the sub-problems were genuinely wasted and worth almost nothing. Where the
count of sub-problems and the cost of the search come apart, it is the count
that is the easier number to improve and the weaker one to quote.

\section{Two obstacles specific to a fractional objective}

Criterion-space search is a known family for optimizing a \emph{linear}
function over the efficient set. It does not transfer to $(P_E)$ by
inspection, for two reasons.

\begin{remark}[$\Phi$ is a function of $x$, not of $\Z(x)$]
\label{rem:notZ}
Two feasible points with the same criterion vector may have different values of
$\Phi$. Enumerating non-dominated vectors and evaluating $\Phi$ on a
representative of each is therefore wrong. Inside every box a genuine
fractional integer program must be solved, and before cutting on an efficient
centre $\tilde{x}$ the sub-problem
$Q(\tilde{x}) = \max \{\, \Phi(x) : x \in \Dset,\ \Z(x) = \Z(\tilde{x}) \,\}$
must be solved first, since the whole slice $\{\Z = \Z(\tilde{x})\}$ is
efficient and is about to be removed.
\end{remark}

\begin{remark}[Box boundaries must be integer-valued]
As argued above, numeric bounds on $\Z_k$ are rational in the fractional case.
Writing the boundaries through $e_k$ restores an exact integral threshold and,
as a by-product, makes fractional criteria work with no special case.
\end{remark}

\section{The algorithm}

\begin{algorithm}[h]
\caption{Criterion-space search for $(P_E)$, optionally seeded}
\label{alg:box}
\begin{algorithmic}[1]
\State $L \gets \{\, \R^p \,\}$
\State $(x_{\mathrm{opt}}, \Phi_{\mathrm{opt}}) \gets$ a \emph{verified
       efficient} seed and its value, or $(\varnothing, -\infty)$
\While{$L \neq \varnothing$}
  \State remove from $L$ a box $B$ of largest inherited bound
  \State solve \eqref{eq:sub} on $B$ with cutoff $\Phi_{\mathrm{opt}}$
  \If{infeasible, or its value $\leqslant \Phi_{\mathrm{opt}}$}
     \State \textbf{continue} \Comment{$B$ is discarded whole}
  \EndIf
  \State let $x_B$ be its maximizer
  \If{$x_B$ is efficient}
     \State $\Phi_{\mathrm{opt}} \gets \Phi(x_B)$; \
            $x_{\mathrm{opt}} \gets x_B$; \ $a \gets x_B$
     \Statex \Comment{$x_B$ already maximizes $\Phi$ on its whole slice, so no $Q$ is owed}
  \Else
     \State $a \gets$ an efficient point dominating $x_B$
     \State solve $Q(a)$ and update $(\Phi_{\mathrm{opt}}, x_{\mathrm{opt}})$
  \EndIf
  \State add to $L$ the $p$ children of $B$ given by \eqref{eq:child}
\EndWhile
\State \Return $(x_{\mathrm{opt}}, \Phi_{\mathrm{opt}})$, proved optimal
\end{algorithmic}
\end{algorithm}

Termination: every pass removes at least the criterion vector $\Z(a)$, and the
non-dominated set of a bounded integer program is finite.

Line 2 is the whole of the hybridization interface. Everything in
Sections~\ref{sec:meta} and~\ref{sec:tcheb} is a different way of filling it,
and the correctness requirement is stated once, in Proposition~\ref{prop:seed}
below: the seed must be \emph{efficient}, not merely good.

\begin{proposition}[Soundness of an arbitrary seed]\label{prop:seed}
If the seed $x_0$ in line 2 is efficient, Algorithm~\ref{alg:box} returns the
same $(x_{\mathrm{opt}}, \Phi_{\mathrm{opt}})$ as the unseeded run, and
$\Phi_{\mathrm{opt}}$ is the global optimum of $(P_E)$. If $x_0$ is merely
feasible, the algorithm may return a strictly suboptimal value.
\end{proposition}

\begin{proof}
$\Phi_{\mathrm{opt}}$ is used only as a cutoff, and a box is discarded when its
bound fails to \emph{beat} it. If $x_0 \in \E$ then
$\Phi(x_0) \leqslant \max_{x \in \E}\Phi(x)$, so no box containing an optimal
criterion vector is ever discarded, and the incumbent is a feasible value of
$(P_E)$ throughout; the two runs therefore agree. If $x_0 \notin \E$ then
$\Phi(x_0)$ may exceed the optimum of $(P_E)$ --- a dominated point can score
arbitrarily well on $\Phi$ --- and the box containing the true optimum is
discarded by the very first bound comparison.
\end{proof}

\noindent
This is why every seed in this note is \emph{verified} before it is used, and
why that verification, not the heuristic, turns out to be the cost of
Section~\ref{sec:meta}.

\section{Computational results: decision space versus criterion space}

All figures come from an exact rational implementation with no floating-point
arithmetic anywhere, so every comparison below is between two exactly solved
problems. Both methods were run on the same instances in the same process.

\begin{table}[h]
\centering
\caption{Decision space versus criterion space, $p = 3$.}
\begin{tabular}{lrrr}
\toprule
instance & decision space & criterion space & ratio \\
\midrule
35 instances, total & $40.65$ s & $8.06$ s & $5.05\times$ \\
heaviest $n=10$     & $23.47$ s & $2.61$ s & $9.01\times$ \\
$n=8$, hardest seed & $2.97$ s  & $0.45$ s & $6.56\times$ \\
$n=10$, easiest seed& $0.36$ s  & $0.41$ s & $0.86\times$ \\
\bottomrule
\end{tabular}
\end{table}

The box search is faster on 34 of the 35, and the single regression is the
instance the decision-space method settles in three iterations, where paying 13
boxes is a net loss. Every answer agrees with the decision-space method and
with an independent scan, and both are proved optimal everywhere.

\begin{table}[h]
\centering
\caption{The margin against the number of criteria, $n = 6$, four instances
each. All twenty agreed and were proved optimal on both sides.}
\begin{tabular}{crrr}
\toprule
$p$ & decision space & criterion space & ratio \\
\midrule
2 & $0.61$ s   & $0.35$ s & $1.74\times$ \\
3 & $1.02$ s   & $0.34$ s & $2.99\times$ \\
4 & $4.21$ s   & $0.61$ s & $6.92\times$ \\
5 & $112.77$ s & $1.09$ s & $103\times$ \\
6 & $26.15$ s  & $1.08$ s & $24.3\times$ \\
\bottomrule
\end{tabular}
\end{table}

Two forces grow with $p$ and pull against each other: the split produces $p$
children, so the list branches wider; but the non-dominated set grows too, so
the decision-space method needs more cuts, and each cut costs it $p$ binaries
and $p+1$ rows. The measurement settles which dominates, and it settles it
against the prediction we made before running it. Branching multiplies a
\emph{count} of cheap, fixed-size models; the binaries multiply the \emph{cost
of one model}, and they accumulate. The worst case observed for the published
method was $p = 5$ on six variables: $110.95$ s against $0.64$ s, a factor of
$173$.

\paragraph{On the magnitudes.} The direction is solid: the per-$p$ median
ratios are $1.35$, $3.2$, $2.1$, $12.6$, $15.8$, and every instance agreed and
was proved optimal. The totals are not stable estimates --- four instances per
$p$ is few, and $110.95$ of the $112.77$ seconds at $p=5$ come from a single
instance. The tables report what was measured, not expected values.

\section{First hybridization: exact--exact}

Since the published method closes in one iteration whenever the maximizer of
$\Phi$ over $\Dset$ happens to be efficient, a hybrid suggests itself: run its
loop, and if it has not closed after $s$ iterations, hand the unfinished work
to the box search. The handover is not a restart. Replaying each cut as a box
subtraction (Proposition~\ref{prop:split} applied to every box of the current
list) reproduces the region the first phase reached, and the incumbent, being
attained at a known efficient point, is a valid seed by
Proposition~\ref{prop:seed}.

Measured on 24 instances against the better of the two pure methods on each:

\begin{table}[h]
\centering
\begin{tabular}{lrr}
\toprule
 & total & within 15\% of the better \\
\midrule
$s = 3$              & $13.14$ s  & 4 of 24 \\
$s = 1$              & $8.47$ s   & 23 of 24 \\
pure box search      & $8.11$ s   & --- \\
the published method & $268.61$ s & --- \\
\bottomrule
\end{tabular}
\end{table}

\noindent
We report this as a negative result. At $s = 1$ the hybrid \emph{ties} the pure
box search --- $8.47$ s against $8.11$ s here, and $1.47$ s against $1.54$ s on
a family built so that $\Phi$ is easy; two differences of about 4\% pointing
opposite ways. The criterion-space search is what matters; this hybrid earns
its place only as insurance for the one-iteration case, not as a third method
that beats both. The first default we chose, $s = 3$, was measurably wrong and
is reported here because it was: $13.14$ s against $8.47$ s.

\section{Second hybridization: exact--metaheuristic}\label{sec:meta}

By Proposition~\ref{prop:seed} the box search will accept any efficient point
as a seed, and by Section~\ref{sec:bound} a good one deletes subtrees. That
makes a metaheuristic the natural first phase --- it is cheap and it is guided
by $\Phi$, which is exactly the quantity the cutoff compares against.

\subsection{The heuristic}

A Pareto local search maintains an archive of mutually non-dominated points.
From several random maximal starting points it walks the neighbourhood of each
archive member, admitting any neighbour the archive does not dominate. The
archive \emph{is} the search: restarting from fresh random points between walks
throws it away, and four times the effort did not compensate for that
(obtaining $0.360$, $0.280$ and $0.458$ of the optimum on three instances where
the archive-driven version obtained the optimum).

The neighbourhood needs \emph{swap} moves ($-1$ on one coordinate, $+1$ on
another), not only unit steps. From a maximal point a unit step up is
infeasible, and a unit step down lowers every criterion with non-negative
coefficients, so the archive refuses it and the walk dies immediately: three
archive points and an incumbent at 13\% of the optimum. With swaps the archive
holds 28--69 points and the incumbent reaches 60--100\%.

\subsection{Where correctness lives}

The archive is filtered by dominance \emph{among the points it has seen}, which
is not the exact efficiency test. Proposition~\ref{prop:seed} makes that gap
fatal if ignored, so nothing crosses unverified: the best archive member is put
through the exact test, and if it fails, it is walked to an efficient point
dominating it and that point's value is used instead. The heuristic decides
\emph{where to look}; it never decides what is true.

The gap is not hypothetical. Over a wide sweep, one archive member in roughly
2000 turned out not to be efficient --- reproducibly, a concrete case being a
four-variable instance whose archive retains $(2,1,1,1)$. Rare is not never.

\subsection{What it costs, and what it is worth}

The cost is \emph{not} the walk, which takes $0.01$--$0.08$\,s and is barely
sensitive to its budget. We assumed the opposite and optimized the walk first;
the timings then showed the total \emph{falling as the walk's budget rose},
because a poor archive yields candidates needing longer repair chains, each one
an integer program. The cost is the verification --- the price of
Proposition~\ref{prop:seed}. Examining the single best archive member gave
exactly the same incumbent as examining four on every instance tried, at a
quarter to a half of the cost.

\begin{table}[h]
\centering
\caption{Exact--metaheuristic hybrid, 18 instances. Same optimum everywhere,
all proved optimal.}
\begin{tabular}{lrrr}
\toprule
instance & exact only & metaheuristic $+$ exact & ratio \\
\midrule
total, 18 instances & $9.50$ s & $7.60$ s & $1.25\times$ \\
\midrule
hardest $n=10$ & 49 boxes, $3.02$ s & 34 boxes, $1.94$ s & $1.56\times$ \\
$n=8$, seed 0  & 16 boxes, $0.53$ s & 10 boxes, $0.34$ s & $1.53\times$ \\
$p=7$, seed 1  & 197 boxes, $1.36$ s & 120 boxes, $0.95$ s & $1.43\times$ \\
$p=5$, seed 1  & 6 boxes, $0.02$ s  & 6 boxes, $0.03$ s   & $0.56\times$ \\
\bottomrule
\end{tabular}
\end{table}

It wins where the search is expensive and loses where there was nothing to win,
which is what a nearly fixed first-phase cost predicts. The \emph{ceiling},
measured by handing the search the true optimum for free, is $1.42\times$; at
$1.25\times$ the heuristic realizes 88\% of it.

\section{The asymmetry: a seed is worth what the exact method can do with it}
\label{sec:asymmetry}

The ceiling experiment can be run on the published method too, and this is the
comparison we consider the most transferable result in this note.

\begin{table}[h]
\centering
\caption{The \emph{same} free optimum, handed to each exact method at
iteration 1.}
\begin{tabular}{lrr}
\toprule
exact method & effect of a free optimal incumbent & mechanism \\
\midrule
criterion-space box search & $1.42\times$, up to 52\% fewer boxes
  & boxes discarded whole \\
decision-space cut method  & \textbf{zero} iterations saved
  & none \\
\bottomrule
\end{tabular}
\end{table}

\noindent
The reason is structural. What the decision-space method cuts at iteration $l$
is decided by the efficiency test applied to $x_l$, through
\eqref{eq:cut}; $\Phi_{\mathrm{opt}}$ enters the sub-problem only as a cutoff
inside a branch-and-bound that would have found the same $x_l$ anyway. A box,
by contrast, is compared against the incumbent \emph{before} it is ever split,
and failing that comparison deletes it together with every descendant it would
have had.

So ``hybridize an exact method with a metaheuristic'' is not one question with
one answer. The heuristic is identical in both rows above; what differs is
whether the exact method has a mechanism through which a lower bound can remove
work. Before building such a hybrid, the mechanism is what should be checked.

\subsection{Two ideas that feed the same mechanism do not add}
\label{sec:substitutes}

The same reasoning has a consequence that is easy to miss when improvements are
collected one at a time. The early exit of Section~\ref{sec:exit} and the
metaheuristic seed of Section~\ref{sec:meta} look independent --- one prunes at
the top of the list, the other starts the incumbent higher --- but they act
through the \emph{same} mechanism, the bound comparison of
Section~\ref{sec:bound}. So they should not be expected to compose, and
measured, they do not:

\begin{table}[h]
\centering
\caption{Sub-problems solved over the 18 instances, with each idea alone and
with both.}
\begin{tabular}{lrr}
\toprule
 & boxes solved & reduction \\
\midrule
neither                      & 733 & --- \\
early exit only              & 549 & 25\% \\
metaheuristic seed only      & 562 & 23\% \\
both                         & 539 & 26\% \\
\bottomrule
\end{tabular}
\end{table}

\noindent
Either alone removes about a quarter of the sub-problems; together they remove
about a quarter. On top of the seed, the exit is worth 4\% rather than 25\%,
because a good incumbent arrives early enough that the hopeless tail barely
forms --- there is nothing left to skip. They are \textbf{substitutes, not
complements}.

We record this because the arithmetic of stacked improvements is where
plausible reasoning goes wrong most cheaply. Two ideas add only when they feed
different mechanisms; two that feed the same one are one idea, measured twice.

\section{Third hybridization: the Tchebychev reference}\label{sec:tcheb}

The surrounding literature generates efficient points with the augmented
weighted Tchebychev program, so a fair evaluation of the two hybrids above has
to be made against it and not only against a weighted sum. We therefore
implement it and measure it, rather than cite it.

\subsection{The program}

With $z^{*}$ the ideal point, $z^{*}_k = \max\{\Z_k(x) : x \in \Dset\}$,
weights $w > 0$ and an augmentation $\rho > 0$:
\begin{equation}\label{eq:tcheb}
  \min_{x \in \Dset} \quad \max_{k} \; w_k\bigl(z^{*}_k - \Z_k(x)\bigr)
  \;+\; \rho \sum_{k=1}^{p} \bigl(z^{*}_k - \Z_k(x)\bigr),
\end{equation}
which linearizes with a single continuous variable $\lambda$ subject to
$\lambda \geqslant w_k (z^{*}_k - \Z_k(x))$ for each $k$. Each weight vector
therefore costs one integer program on a model with $n+1$ variables and $p$
extra rows, plus $p$ programs paid once for $z^{*}$.

\begin{proposition}\label{prop:tcheb}
For any $w > 0$ and any $\rho > 0$, every optimal solution of
\eqref{eq:tcheb} is efficient.
\end{proposition}

\begin{proof}
Let $x^{*}$ be optimal and suppose some $y \in \Dset$ dominated it, so
$\Z(y) \geqslant \Z(x^{*})$ with at least one strict inequality. Then
$z^{*}_k - \Z_k(y) \leqslant z^{*}_k - \Z_k(x^{*})$ for every $k$, so the
$\max$ term does not increase; and since $\rho > 0$ and at least one
inequality is strict, the sum term \emph{strictly} decreases. The objective at
$y$ would be strictly smaller than at $x^{*}$, contradicting optimality.
\end{proof}

\noindent
The augmentation is exactly what makes the second half of that argument
available: with $\rho = 0$ the objective at $y$ need only be no larger, and a
merely \emph{weakly} efficient point can minimize the program. Our
implementation therefore refuses $\rho \leqslant 0$ and refuses a zero weight
rather than returning a point whose status is in doubt. It also refuses
fractional criteria, for which $z^{*}_k - \Z_k(x)$ is not linear and the rows
above are not constraints of an integer linear program; the references
assume linear criteria as well.

\subsection{As a generator it is strictly stronger than a weighted sum}

A weighted sum maximizes only at supported efficient points. Varying $w$ in
\eqref{eq:tcheb} reaches unsupported ones too, and that is what the literature
uses it for.

\begin{table}[h]
\centering
\caption{Reach from the same grid of 27 weight vectors, 20 random instances.}
\begin{tabular}{lrr}
\toprule
 & efficient points reached & of which unsupported \\
\midrule
weighted sum         & $66 / 175$  & $\mathbf{0 / 54}$ \\
augmented Tchebychev & $112 / 175$ & $29 / 54$ \\
\bottomrule
\end{tabular}
\end{table}

\noindent
All 540 optima returned were efficient, as Proposition~\ref{prop:tcheb}
requires. A wider sweep over 40 instances agrees: $216/369$ against $135/369$,
56 unsupported against 0, and 1080 optima with no inefficient one. On the
reference instance of the published paper, three of the seven efficient points
are unsupported; the Tchebychev program reaches all three and no strictly
positive weighted sum reaches any.

\subsection{As a seed it loses to not seeding at all}

\begin{table}[h]
\centering
\caption{Seeding the box search, 18 instances. ``Exactly optimal'' counts the
instances where the seed already attains $\Phi^{*}$.}
\begin{tabular}{lrrr}
\toprule
seed & total & seed cost & exactly optimal \\
\midrule
none (pure box search) & $9.64$ s  & ---      & --- \\
Pareto local search    & $\mathbf{7.97}$ s $(1.21\times)$ & $0.57$ s & 12 of 18 \\
augmented Tchebychev   & $13.50$ s $(0.71\times)$ & $4.25$ s & 5 of 18 \\
\bottomrule
\end{tabular}
\end{table}

\noindent
Both halves of the loss are structural rather than artefacts of the
implementation.

\paragraph{Cost.} It is $p+1$ integer programs on an enlarged model. It is
\emph{not} the reference point: computing $z^{*}$ took $0.01$--$0.13$\,s of the
seed against $0.13$--$0.85$\,s for the scalarizations themselves, so replacing
$z^{*}$ by a cheaper reference would not change the verdict.

\paragraph{Quality.} This is the real finding. Program~\eqref{eq:tcheb} steers
by criterion-space geometry relative to $z^{*}$ and \textbf{never looks at
$\Phi$}. The local search of Section~\ref{sec:meta} is guided by $\Phi$ at
every move. A good spread of efficient points is not a good $\Phi$, and by
Remark~\ref{rem:notZ} it cannot be made into one by post-processing: $\Phi$ is
not a function of $\Z(x)$, so two points the generator regards as
interchangeable may differ arbitrarily in the only quantity the search cares
about.

We therefore keep the Tchebychev program as the reference the evaluation needs,
and do not wire it into the default hybrid.

\subsection{The two generators compared directly}

\begin{table}[h]
\centering
\caption{As generators of a subset of $\E$, 20 instances at $n = 4$.}
\begin{tabular}{lrrl}
\toprule
 & coverage of $\E$ & time & certification \\
\midrule
augmented Tchebychev & 64\% & $2.1$ s  & every point efficient by construction \\
Pareto archive       & 99\% & $0.04$ s & $\approx 1$ member in 2000 not efficient \\
\bottomrule
\end{tabular}
\end{table}

\noindent
Cheap, near-complete and uncertified against expensive, partial and certified.
That is the trade the exact--metaheuristic hybrid closes: it takes the cheap,
near-complete generator and buys the missing certificate with a single exact
test, which is why the verification and not the walk is its dominant cost.

This also delimits the expected result of the thesis --- ``a subset of
efficient solutions generated by hybrid algorithms''. Both columns produce such
a subset, and they are not interchangeable: one is a certified partial sample
of the front, the other a near-complete sample whose membership is a
conjecture until tested.

\section{Two measurement errors, and what they cost}

We report these because the direction of a result is worth less than the
knowledge of how it was checked.

\paragraph{An unsound quality metric.} Seed quality was first scored by the
ratio $\Phi(\text{seed}) / \Phi^{*}$. That ratio \emph{inverts when $\Phi^{*}$
is negative}, which it is on several instances in these families: a strictly
worse seed can score above~1. All seed-quality figures in this note use the
non-negative gap $(\Phi^{*} - \Phi(\text{seed}))/|\Phi^{*}|$, which is
well defined whatever the sign and is zero exactly when the seed is optimal.
No published number changed direction under the correction, but the metric was
wrong and one table read $1.286$ for a seed that was in fact 29\% below the
optimum.

\paragraph{Seed quality is not seed value.} On a looser instance family, the
Pareto seed is \emph{exactly optimal on 15 of 24 instances} and still buys only
$1.05\times$, because that search is not bound-limited: there are few boxes
whose bound the incumbent can defeat. Quality is a property of the heuristic;
value is a property of the pair. What is stable across all families we tried is
the \emph{ordering} --- Tchebychev loses to no seed at all ($0.73\times$,
$0.91\times$), the Pareto seed ties or wins --- not the size of any win. Single
speed-up figures quoted without their instance family should be read with this
in mind, including ours.

\section{Summary of what was and was not gained}

\begin{table}[h]
\centering
\caption{Every hybridization attempted, with its verdict.}
\begin{tabular}{lll}
\toprule
hybridization & verdict & effect \\
\midrule
decision space $\to$ criterion space & \textbf{win} & $5.05\times$, up to $173\times$ \\
exact--exact handover ($s=1$)        & tie          & $8.47$ s vs $8.11$ s \\
exact--exact handover ($s=3$)        & \textbf{loss} & $13.14$ s vs $8.47$ s \\
exact--metaheuristic (Pareto archive) & \textbf{win} & $1.25\times$ of a $1.42\times$ ceiling \\
exact--Tchebychev seed               & \textbf{loss} & $0.71\times$ \\
Tchebychev as a generator            & \textbf{win} & reaches unsupported points \\
incumbent seeding in decision space  & \textbf{loss} & zero iterations saved \\
early exit on a hopeless bound       & partial      & 25\% fewer boxes, no time \\
early exit \emph{on top of} the seed  & \textbf{loss} & 4\% fewer boxes: a substitute \\
\bottomrule
\end{tabular}
\end{table}

\noindent
Of the nine, three are wins, four are losses, one is a tie and one pays in
sub-problems but not in time. They are reported at the same length as the wins
because they carry the same information: a hybridization pays only where the
exact method has a mechanism the first phase can feed, that mechanism is
cheaper to identify than the hybrid is to build and measure, and two ideas
feeding the same mechanism are one idea counted twice.

\section{Prior work, and what remains to be checked}

Criterion-space search for optimizing a function over the efficient set of a
multi-objective integer program is an established family, and the surveys
report it as currently more effective than decision-space search. Those methods
are for \emph{linear} objectives. The literature on a linear \emph{fractional}
objective over an integer efficient set --- including the branch-and-cut
methods this note compares against, and the extensions to fractional criteria
and to quadratic objectives --- works in decision space. The augmented weighted
Tchebychev program of Section~\ref{sec:tcheb} is standard and is used here
exactly as the references use it; nothing about it is claimed as new. The
combination treated in Section~\ref{sec:box} appears not to have been
addressed, for the reason given in Remark~\ref{rem:notZ}: $\Phi$ is a function
of $x$ and not of $\Z(x)$, so the transfer is not immediate.

\textbf{This is a statement about what we did not find, not a claim of
novelty.} A proper bibliographic review, including the sources behind
paywalls, remains to be done before any such claim is made.

\section{Reproducibility}

The implementation is pure Python with no third-party dependency, computes in
exact rational arithmetic throughout, and carries 55 tests. Several check
properties rather than answers, since a method that silently lost efficient
points would still agree with the optimum on favourable instances:

\begin{itemize}
\itemsep2pt
\item Proposition~\ref{prop:split} directly --- for every centre and every
      feasible point, the point is either removed by the cut and in no child,
      or kept and in exactly one.
\item Proposition~\ref{prop:seed} operationally --- the metaheuristic is
      checked, over 14 instances and against an independent scan, never to hand
      over an inefficient point.
\item Proposition~\ref{prop:tcheb} --- every optimum of \eqref{eq:tcheb} is
      checked efficient against exhaustive enumeration, and the refusals
      ($\rho \leqslant 0$, a zero weight, fractional criteria) are checked to
      fire.
\item The unsupported-reach claim --- on the reference instance, the three
      unsupported efficient points are all reached by \eqref{eq:tcheb} while no
      positive weighted sum reaches any.
\end{itemize}

\noindent
Every comparison in this note is reproduced by scripts in the repository, run
on every commit by continuous integration, which assert that each seeded run
returns the same proved optimum as the unseeded one.

\end{document}
